\clearpage
\section[Rev. 1.41 -- Watchdog esterno STWD100 (U403)]{Rev. 1.41 \textendash{} Watchdog esterno STWD100 (U403) sul reset dell'MCU}

\subsection{Data e revisione}
03/11/2026, revisione 1.41.

\subsection{Titolo}
Watchdog esterno STWD100 (U403) sul reset dell'MCU.

\subsection{Componenti, net e fogli coinvolti}
\begin{itemize}
\item Foglio MCU (\code{MCU.kicad\_sch}) e libreria di progetto \code{UHF\_LIB/SYMBOL/IC.kicad\_sym} (nuovo simbolo \code{IC:STWD100NYWY3F}). Copie precedenti in \code{report\_src/bak\_pre141/} (\code{MCU.kicad\_sch}, \code{IC.kicad\_sym}).
\item Reti: \code{/IF/DBG.nRST} (NRST dell'STM32 U301 pin 7, pull-up R2 \SI{10}{\kilo\ohm}, C3 \SI{100}{\nano\farad}, debugger TC2030 J1 pin 3), nuova rete \code{/MCU/WDT\_WDI}, \code{3V3}, \code{GND}, nuova rete \code{Net-(U403-\textasciitilde{}EN)}.
\item Simbolo nuovo con riquadro \textquotedblleft{}WATCHDOG\textquotedblright{} sotto il blocco FRAM.
\end{itemize}

\subsection{Cosa è cambiato}
Aggiunti tre componenti.

{\small\arrayrulecolor{rulec}
\begin{tabularx}{\linewidth}{@{}>{\hsize=0.3\hsize\raggedright\arraybackslash}X>{\hsize=1.0\hsize\raggedright\arraybackslash}X>{\hsize=0.5\hsize\raggedright\arraybackslash}X>{\hsize=1.9\hsize\raggedright\arraybackslash}X>{\hsize=1.3\hsize\raggedright\arraybackslash}X@{}}
\toprule
\textbf{Rif.} & \textbf{Parte} & \textbf{JLCPCB} & \textbf{Footprint} & \textbf{Note}\\
\midrule
U403 & STMicroelectronics \code{STWD100NYWY3F} & C46043 & \code{IC:SOT95P280X145-5N} & stock JLC 21\,124 al 03/11/2026; stesso SOT23-5 di IC205, IC1, IC5\\
C425 & \SI{100}{\nano\farad} 0402 & C1525 & \code{Capacitor\_SMD:C\_0402\_1005Metric} & disaccoppiamento VCC di U403\\
R413 & \SI{0}{\ohm} 0402 & C17168 & \code{Resistor\_SMD:R\_0402\_1005Metric} & DNP (Value \code{0R} + flag DNP)\\
\bottomrule
\end{tabularx}}

\medskip
Collegamenti:

{\small\arrayrulecolor{rulec}
\begin{tabular}{@{}clll@{}}
\toprule
\textbf{Pin U403} & \textbf{Nome} & \textbf{Collegato a} & \textbf{Rete}\\
\midrule
1 & \textasciitilde{}WDO (open drain, attivo basso) & NRST di U301 (pin 7) & \code{/IF/DBG.nRST}\\
2 & GND & GND & \code{GND}\\
3 & \textasciitilde{}EN & R413 (DNP) $\rightarrow$ 3V3 & \code{Net-(U403-\textasciitilde{}EN)}\\
4 & WDI & U301 PC6 (pin 37) & \code{/MCU/WDT\_WDI}\\
5 & VCC & 3V3, con C425 verso GND & \code{3V3}\\
\bottomrule
\end{tabular}}

\medskip
U301 PC6 (pin 37) era libero con no-connect, che è stato rimosso.

\begin{center}
\begin{adjustbox}{max width=\linewidth}\begin{circuitikz}[scale=0.9, transform shape]
\draw (0,0) rectangle (3,3);
\node at (1.5,1.5) {U403 STWD100};
\node[anchor=west,font=\scriptsize] at (0.05,2.2) {3 EN};
\node[anchor=west,font=\scriptsize] at (0.05,0.8) {4 WDI};
\node[anchor=east,font=\scriptsize] at (2.95,1) {1 WDO};
\node[anchor=north,font=\scriptsize] at (1.5,2.95) {5 VCC};
\node[anchor=south,font=\scriptsize] at (1.5,0.05) {2 GND};
% EN con R413 DNP a 3V3
\draw (0,2.2) -- (-1.5,2.2) to[R, l_=R413 (DNP)] (-1.5,3.8) -- (5.5,3.8);
\node[font=\scriptsize,anchor=south] at (-1.5,3.8) {3V3};
% VCC e C425
\draw (1.5,3) -- (1.5,3.8) node[circ]{};
\draw (4,3.8) node[circ]{} to[C, l_=C425] (4,2.2) node[ground]{};
% GND
\draw (1.5,0) -- (1.5,-0.3) node[ground]{};
% WDI
\draw (0,0.8) -- (-3,0.8) node[anchor=east,font=\scriptsize]{WDT\_WDI $\leftarrow$ PC6 (U301)};
% WDO e pull-up R2
\draw (3,1) -- (5.5,1) node[circ]{} -- (7,1) node[anchor=west,font=\scriptsize]{nRST};
\draw (5.5,3.8) to[R, l=R2] (5.5,1);
\end{circuitikz}\end{adjustbox}
\end{center}

\subsection{Motivazione}
Brik81 vuole un watchdog esterno indipendente dall'MCU. Varianti JLC confrontate: \code{C46043} STWD100NYWY3F (\SI{1,6}{\second}, 21\,124 pz) scelto; \code{C1852782} STWD100NXWY3F (\SI{102}{\milli\second}, 1\,206 pz), \code{C1354132} STWD100YNXWY3F (\SI{102}{\milli\second} automotive, 209 pz) e \code{C2842415} STWD100NPWY3F (\SI{3,4}{\milli\second}, 88 pz) scartati: timeout troppo corti per scritture flash e trasmissioni.

Il pin EN non è collegato a un GPIO per evitare che un firmware fuori controllo disabiliti il proprio watchdog; per debug e bootloader UART si monta R413 (EN alto = watchdog disabilitato). Vincolo di Brik81: passivi obbligatoriamente 0402.

Dati dal datasheet ST STWD100, DocID14134 Rev 10 (settembre 2016):
\begin{itemize}
\item Figure 13, p.\,18, codice d'ordine: N = uscita open drain attiva bassa; Y = $t_{WD}$ \SI{1,6}{\second} e $t_{PW}$ \SI{210}{\milli\second}; WY = SOT23-5; 3 = \SIrange{-40}{125}{\celsius}.
\item Table 1, p.\,3, pinout: 1 WDO, 2 GND, 3 EN, 4 WDI, 5 VCC.
\item Table 4, p.\,12--13: vedi tabella sotto.
\item §3.2: con uscita open drain serve un pull-up esterno (qui R2 \SI{10}{\kilo\ohm} già presente su NRST).
\item §3.1: WDI va commutato entro $t_{WD}$; sulla versione Y conta anche il fronte di discesa.
\end{itemize}

{\small\arrayrulecolor{rulec}
\begin{tabular}{@{}lrrr@{}}
\toprule
\textbf{Parametro (versione Y)} & \textbf{min} & \textbf{tip} & \textbf{max}\\
\midrule
$t_{WD}$ & \SI{1,12}{\second} & \SI{1,6}{\second} & \SI{2,24}{\second}\\
$t_{PW}$ & \SI{140}{\milli\second} & \SI{210}{\milli\second} & \SI{280}{\milli\second}\\
Pull-down interno EN & \SI{32}{\kilo\ohm} & \SI{63}{\kilo\ohm} & \SI{100}{\kilo\ohm}\\
$V_{CC}$ & \SI{2,7}{\volt} & & \SI{5,5}{\volt}\\
$I_{CC}$ & & \SI{13}{\micro\ampere} & \\
\bottomrule
\end{tabular}}

\medskip
Impulso minimo su WDI ed EN \SI{1}{\micro\second}; glitch $<$\SI{100}{\nano\second} ignorati.

Approvazioni di Brik81 (03/11/2026): \textquotedblleft{}Si va bene mettiamo il C46043, come va connesso? se servono componenti passivi OBBLIGATORIO componenti da 0402 ok?\textquotedblright{} e, sullo schema di collegamento proposto (WDI su PC6, R413 0R DNP su EN, C425 \SI{100}{\nano\farad}), \textquotedblleft{}aggiungilo tu e io verificherò\textquotedblright{}.

Eseguito da Claude Code il 03/11/2026; Brik81 verificherà in KiCad.

\subsection{Impatti e criticità}
\begin{critbox}
Con R413 non montato il watchdog è sempre attivo: in laboratorio la programmazione via SWD con l'MCU fermo e il bootloader UART (BOOT0) provocano un reset ogni circa \SI{1,6}{\second}. Montare R413 durante lo sviluppo e toglierlo per il volo.
\end{critbox}
\begin{notebox}
Dopo un reset dell'MCU PC6 è in alta impedenza: se il firmware non riparte entro \SI{1,12}{\second} il watchdog resetta di nuovo (voluto).

Al reset dell'MCU PC2 e PC4 vanno in alta impedenza: PA spento (pull-up R702), radio sotto il controllo di IC205 (modo A): comportamento sicuro.
\end{notebox}

\textbf{Note per il firmware} (continuano la serie F-01\ldots F-10 del §41):
\begin{notebox}
\textbf{F-11} Impulso su PC6 (\code{WDT\_WDI}) almeno ogni circa \SI{1}{\second}, anche durante operazioni lunghe ($t_{WD}$ minimo garantito \SI{1,12}{\second}).\\
\textbf{F-12} Primo impulso entro \SI{1,12}{\second} dall'accensione, prima delle inizializzazioni lente.\\
\textbf{F-13} All'avvio leggere \code{RCC\_CSR} per la causa del reset (pin NRST = watchdog esterno) e riportarla in telemetria.
\end{notebox}

\subsection{Verifiche fatte}
\begin{itemize}
\item Netlist prima/dopo: unica rete rimossa \code{unconnected-(U301-PC6-Pad37)}; aggiunte U403 pin 1 su \code{/IF/DBG.nRST}, U301 pin 37 e U403 pin 4 su \code{/MCU/WDT\_WDI}, C425, U403 pin 5 e R413 pin 2 su \code{3V3}, C425 pin 2 e U403 pin 2 su \code{GND}, R413 pin 1 e U403 pin 3 su \code{Net-(U403-\textasciitilde{}EN)}; nient'altro cambia.
\item Schema stampato in PDF e controllato a vista.
\item ERC: solo avvisi già presenti nel progetto (\code{lib\_symbol\_mismatch} di C\_Small, \code{pin\_to\_pin} con J1).
\end{itemize}

\subsection{Verifiche da fare}
\begin{itemize}
\item Misura M-04 in collaudo: con R413 DNP e firmware bloccato, oscilloscopio su NRST: impulso basso di circa \SI{210}{\milli\second} (\SIrange{140}{280}{\milli\second}) ogni circa \SI{1,6}{\second} (\SIrange{1,12}{2,24}{\second}); con R413 montato nessun reset.
\item Verifica di Brik81 in KiCad.
\end{itemize}
